% =============================================================================
\section{Approche hybride}
% =============================================================================

% Hybridation : exploration + intensification.
\begin{frame}{Hybridation : combiner exploration et intensification}
  \footnotesize
  \begin{columns}[T,onlytextwidth]
    \begin{column}{0.52\textwidth}
      \begin{block}{Hybrides de métaheuristiques}
        \begin{itemize}
          \item \textbf{Tabou $\times$ Recuit} : voisinage tabou + acceptation
            de Metropolis.
          \item \textbf{Génétique $\times$ Tabou} : recherche tabou courte sur
            le meilleur individu (stratégie lamarckienne).
          \item \textbf{Génétique $\times$ Recuit} : recuit court sur le
            meilleur individu.
          \item \textbf{Fourmis $\times$ Tabou} : intensification de la
            meilleure fourmi.
        \end{itemize}
      \end{block}
    \end{column}
    \begin{column}{0.46\textwidth}
      \begin{exampleblock}{Principe}
        \footnotesize
        \begin{itemize}
          \item une méthode \textbf{globale} explore largement (AG, ACO) ;
          \item une méthode \textbf{locale} intensifie les bonnes zones
            (Tabou, Recuit) ;
          \item l'intensification corrige la faible pression de sélection du
            génétique.
        \end{itemize}
      \end{exampleblock}
      \begin{alertblock}{Apport mesuré}
        \footnotesize Génétique $\times$ Tabou : écart du génétique divisé par
        $\sim$$10^4$ ; Fourmis $\times$ Tabou : écart ACO divisé par
        $\sim$1,7.
      \end{alertblock}
    \end{column}
  \end{columns}
  \source{Distinct de l'hybride d'origine Tabou $\times$ Recuit ; les trois
    autres hybrides ont été construits pour ce projet.}
\end{frame}

% SMA Mesa : architecture de coopération.
\begin{frame}{Système multi-agents : tableau noir et coordination}
  \footnotesize
  \begin{columns}[T,onlytextwidth]
    \begin{column}{0.52\textwidth}
      \centering
      \begin{tikzpicture}[
        ag/.style={draw=mcmExample,rounded corners=2pt,fill=mcmExample!12,
                   minimum height=0.68cm,minimum width=2.35cm,align=center,font=\scriptsize},
        bb/.style={draw=mcmPrimary,rounded corners=2pt,fill=mcmPrimaryMuted,
                   minimum height=1.45cm,minimum width=2.7cm,align=center,font=\scriptsize,thick},
        co/.style={draw=mcmAlert,rounded corners=2pt,fill=mcmAlert!12,
                   minimum height=0.75cm,minimum width=2.7cm,align=center,font=\scriptsize,thick},
        flow/.style={-{Latex},thick,draw=mcmPrimary}]
        \node[ag] (a1) {Agent recuit};
        \node[ag,below=0.28cm of a1] (a2) {Agent tabou};
        \node[ag,below=0.28cm of a2] (a3) {Agent génétique};
        \node[bb,right=1.0cm of a2] (bb) {\textbf{Tableau noir}\\[2pt]meilleure\\solution versionnée};
        \node[co,above=0.35cm of bb] (co) {\textbf{Coordinateur}};
        \draw[flow] (a1.east) -- (bb.west|-a1);
        \draw[flow] (a2) -- (bb);
        \draw[flow] (a3.east) -- (bb.west|-a3);
        \draw[flow,draw=mcmAlert] (co) -- (bb);
      \end{tikzpicture}
    \end{column}
    \begin{column}{0.46\textwidth}
      \begin{block}{Les briques}
        \footnotesize
        \begin{itemize}
          \item \textbf{Tableau noir} : meilleure solution partagée,
            versionnée.
          \item \textbf{Agents} : une stratégie chacun, travaillant par
            \textbf{quanta}.
          \item \textbf{Coordinateur} : migration du meilleur ; redémarrage
            si la version stagne.
        \end{itemize}
      \end{block}
      \begin{exampleblock}{Déterminisme}
        \footnotesize Générateur aléatoire propre à chaque agent : résultats
        reproductibles par graine.
      \end{exampleblock}
    \end{column}
  \end{columns}
  \source{Implémentation Mesa 3.5 : deux SMA (agents simples et agents
    métaheuristiques) coopérant par tableau noir.}
\end{frame}

% Benchmark grande échelle.
\begin{frame}{Benchmark à budget égal : 300 patients, 80 vacations}
  \small
  \begin{columns}[T,onlytextwidth]
    \begin{column}{0.54\textwidth}
      \centering
      \renewcommand{\arraystretch}{1.12}
      \footnotesize
      \begin{tabular}{lrrr}
        \toprule
        Méthode & Écart moyen & Rang & IQR \\
        \midrule
        Recuit simulé      & \textbf{0,015} & \textbf{1,25} & 0,008 \\
        Tabou $\times$ Recuit & \textbf{0,018} & 1,75 & 0,013 \\
        Tabou              & 0,240 & 3,55 & 0,070 \\
        Génétique $\times$ Tabou & 0,261 & 3,75 & 0,106 \\
        SMA                & 1,078 & 5,05 & 0,194 \\
        SMA $\times$ Méta  & 1,196 & 6,15 & 0,154 \\
        Génétique $\times$ Recuit & 2,031 & 6,50 & 0,323 \\
        Fourmis $\times$ Tabou & 1\,006,0 & 8,00 & 180,9 \\
        Fourmis (ACO)      & 2\,031,2 & 9,00 & 256,1 \\
        Génétique          & 2\,836,2 & 10,00 & 380,1 \\
        \bottomrule
      \end{tabular}
    \end{column}
    \begin{column}{0.44\textwidth}
      \centering
      \includegraphics[width=\linewidth,height=0.30\textheight,keepaspectratio]{grande_convergence.png}\\[2pt]
      {\tiny\color{mcmPrimary}\textbf{Convergence médiane} (20 graines, 3 s).}\\[4pt]
      \begin{alertblock}{Lecture}
        \footnotesize Le recuit et Tabou $\times$ Recuit dominent : écarts
        $<0{,}02$ sur une référence de $\sim$38\,800. ACO et génétique
        décrochent (forte variabilité).
      \end{alertblock}
    \end{column}
  \end{columns}
  \source{20 graines, budget 3 s, référence $-38\,847{,}195$ ; écart à la
    meilleure solution observée.}
\end{frame}

% Benchmark échelle réelle.
\begin{frame}{Échelle réelle : 582 patients, 280 vacations}
  \small
  \begin{columns}[T,onlytextwidth]
    \begin{column}{0.5\textwidth}
      \renewcommand{\arraystretch}{1.2}
      \footnotesize
      \begin{tabular}{lrr}
        \toprule
        Méthode & Médiane & Rang moyen \\
        \midrule
        Recuit simulé      & $-16{,}3$ & 1,85 \\
        Tabou $\times$ Recuit & $-14{,}9$ & 2,05 \\
        Tabou              & $-16{,}4$ & 2,65 \\
        Génétique $\times$ Tabou & $-23{,}9$ & 3,50 \\
        SMA                & $-91{,}8$ & 5,85 \\
        Fourmis (ACO)      & $-3\,438{,}8$ & 9,70 \\
        \bottomrule
      \end{tabular}
      \\[4pt]
      {\tiny 20 graines ; référence $-8{,}783$.}
    \end{column}
    \begin{column}{0.48\textwidth}
      \centering
      \includegraphics[width=\linewidth,height=0.27\textheight,keepaspectratio]{reelle_convergence.png}\\[2pt]
      \begin{alertblock}{Un couloir faisable étroit}
        \footnotesize La meilleure solution est faisable
        ($\text{violations} = (0,0)$, lits à 42/42). Les méthodes qui
        \textbf{acceptent des dégradations} longent cette frontière ; les
        constructions gloutonnes échouent.
      \end{alertblock}
    \end{column}
  \end{columns}
  \source{Parquet EDA prétraité (40 jours, capacité 600 min/vacation), 4 s par
    méthode et par graine.}
\end{frame}

